\section{Pseudocode}

\begin{algorithm}
[!htb]
\caption{Worker pool setup pseudocode}
\label{pseudo:setup}
\begin{algorithmic}
  [1] \STATE Prepare shared spatial and crime data \STATE Create a pool of worker
  threads with shared data

  \FORALL{tasks to be evaluated} \STATE Package task-specific parameters \STATE
  Submit task to the worker pool \ENDFOR

  \STATE Wait for all workers to complete \STATE Collect results from all
  tasks \STATE Shut down the worker pool
\end{algorithmic}
\end{algorithm}

\begin{algorithm}
[htb!]
\caption{Bruteforce pseudocode}
\label{pseudo:bruteforce}
\begin{algorithmic}
  [1] \STATE Receive task with camera location, distance, and distance weight
  \STATE Access shared data (buildings, bounding box, crimes, crime counts)

  \STATE Generate camera coverage at the given location \STATE Compute the camera
  score \STATE Send result back to the parent process
\end{algorithmic}
\end{algorithm}

\begin{algorithm}
[htb!]
\caption{Hillclimb pseudocode}
\label{pseudo:hillclimb}
\begin{algorithmic}
  [1]

\STATE Pick a random starting point from the grid 
\STATE Generate initial camera coverage and compute score 
\STATE Initialize visited points set and simulation sequence

    \WHILE{step count $<$ MAXSTEPS - 1} 
        \STATE Define all possible movement directions 
        \STATE Evaluate all moves from current position in parallel
        \STATE Filter out moves leading to buildings or previously visited points
            \IF{no valid moves remain} 
                \STATE Stop simulation 
            \ENDIF

        \STATE Sort candidate moves by descending score 
        \STATE Pick the first move that: 
            \STATE \quad increases score OR 
            \STATE \quad keeps score equal but reduces distance traveled

        \IF{no improving move found} 
            \STATE Stop simulation 
        \ENDIF

        \STATE Move to the selected point 
        \STATE Record the score for this step
        \STATE Mark the point as visited 
        \STATE Increment step counter 
    \ENDWHILE

    \STATE Send full simulation sequence back to the parent process
\end{algorithmic}
\end{algorithm}

\begin{algorithm}
[htb!]
\caption{Depth-First search pseudocode}
\label{pseudo:dfs}
\begin{algorithmic}
  [1] \STATE Select a random starting point from the grid \STATE Generate camera
  coverage at starting point \STATE Compute initial camera score

  \STATE Initialize visited set \STATE Initialize simulation sequence \STATE
  Initialize best score with initial score

  \STATE Mark starting point as visited \STATE Add initial score to simulation
  sequence

  \STATE Recursively explore neighboring grid points:

  \STATE \quad Stop exploration if depth $\ge$ MAXSTEPS $-1$ \STATE \quad
  Stop exploration if simulation length $\ge$ MAXSTEPS $-1$

  \STATE \quad From current point, evaluate all movement directions \STATE
  \quad Discard invalid moves and already visited points

  \STATE \quad Sort remaining candidates by: \STATE \qquad descending score
  \STATE \qquad ascending total distance

  \STATE \quad For each candidate move (in sorted order): \STATE \qquad Mark
  point as visited \STATE \qquad Add score to simulation sequence \STATE
  \qquad Update best score if improved \STATE \qquad Recursively continue
  exploration from that point

  \STATE Add best score found to simulation sequence \STATE Send simulation sequence
  to parent process
\end{algorithmic}
\end{algorithm}